\chapter{Equity Track: GAT Relational Factors}
\label{ch:equity}

The equity track asks whether propagating island-alpha information over a
correlation graph produces a factor worth having, and --- the sharper
question, RQ1 of \cref{ch:intro} --- whether \emph{learning} the
propagation weights matters. This
chapter reports the experiments in the order they constrained each other:
the controls that license any positive claim (\cref{sec:eq-controls}),
the first real-data run (\cref{sec:eq-first}), the ablation that
attributed its weaknesses (\cref{sec:eq-ablation}), the seed and device
studies that set the reporting discipline (\cref{sec:eq-variance}),
hyperparameter selection under that discipline
(\cref{sec:eq-hp}), and finally what the trained attention actually does
(\cref{sec:eq-attention}). \Cref{tab:equity-summary} states the headline
up front.

\begin{table}[htbp]
  \centering
  \small
  \caption{Equity track headline results (real yfinance data, 49 names, IC loss).}
  \label{tab:equity-summary}
  \begin{tabular}{l>{\raggedright\arraybackslash}p{0.28\textwidth}>{\raggedright\arraybackslash}p{0.30\textwidth}}
    \toprule
    metric & value & note \\
    \midrule
    Gates passed & 3 / 4 & Uniqueness + Consistency + Robustness pass; Value-added fails \\
    Attention value-add & 30/30 seeds positive & vs the uniform anchor (no learned attention) \\
    Composite best-epoch validation IC & 0.0756 & best-by-valid-IC checkpoint selection \\
    Composite OOS Sharpe & 1.42 & vs best single 2.88 (beats 9 of 10 single factors) \\
    Graph & correlation top-k (~50 names + GICS sectors) & estimated backbone (not physical) \\
    Data & real yfinance 2021-01$\to$2026-06 & 49 names; IC loss; 50 epochs \\
    \bottomrule
  \end{tabular}
\end{table}


\section{Controls before claims}
\label{sec:eq-controls}

No positive number in this chapter means anything unless the pipeline can
be shown to find nothing where nothing exists. Three controls establish
this.

First, the synthetic negative control: fifty random-walk price series run
through the full pipeline --- graph construction, training, gates. The
gates correctly report no edge. Value-added and Consistency fail (OOS IC
around $-0.06$, negative Sharpe value-add), while Uniqueness and
Robustness pass, as they should: they measure structural properties of a
factor, not the presence of signal.

Second, the automated pair of controls that runs in the test suite on
every change. Shuffling each snapshot's labels across nodes and retraining
must collapse the validation IC to zero, and it does, for both loss
functions. Planting a recoverable signal --- the label set equal to a
standardised input feature --- must let the same training loop recover it,
and it does. The second half is what makes the first half meaningful: the
IC collapses on shuffled labels because there is nothing to learn, not
because the trainer is broken.

Third, a convergence probe for the training objective itself. On the
planted-signal panel the IC loss recovers the signal essentially
completely (validation IC 0.99 and above at learning rate $10^{-3}$),
where MSE plateaus at 0.33. The IC loss is unstable early at higher
learning rates --- at $5\times10^{-3}$ it bounces near zero for some forty
epochs before converging --- which fixed the default learning rate at
$10^{-3}$ and the IC loss as the default objective.

\section{First real-data run}
\label{sec:eq-first}

The first run on market data used daily bars for the fifty-name universe
(one name, ORCL, failed to download, leaving 49 names, 1{,}364 trading
days and 66{,}836 panel rows spanning January 2021 to June 2026), the ten
island alphas as node features, a static correlation graph built from
data before the split date, and IC-loss training for fifty epochs.

Training behaved like the textbook picture. The training loss fell
monotonically while the validation IC rose to a peak of 0.0756 at epoch
32 and then decayed to 0.0418 by the final epoch --- plain overfitting.
Best-epoch checkpointing therefore deployed the epoch-32 weights; had the
pipeline returned the final epoch, as an earlier version of the code
silently did, the deployed model would have carried barely half the
validation IC. The checkpoint fix paid for itself on its first contact
with real data.

The four gates split three to one. Consistency (score 0.63), Uniqueness
(maximum absolute correlation against any single alpha, 0.257) and
Robustness (0.68) pass; Value-added fails, with a composite OOS Sharpe of
1.42 against 2.88 for the best single island alpha. The honest reading
cuts both ways. The composite is a real, unique, consistent signal ---
the second-best of the eleven factors on the board, beating nine of the
ten singles it was built from, and correlated with none of them above
0.26. It is genuinely new information rather than a repackaging of its
inputs. But the strictest bar --- beat the maximum of ten competitors ---
stayed open here and, as the rest of this chapter shows, never closed.

The run also posed the question that drove the next experiment: the
validation IC of 0.076 decayed to an OOS IC of 0.007 over the roughly
1.6-year OOS window. Was the frozen correlation graph going stale, or the
frozen model weights?

\section{Attributing the decay: graph versus model}
\label{sec:eq-ablation}

A $2\times2$ ablation crossed the graph axis (static, frozen at the split
date, versus dynamic, rebuilt per snapshot from a trailing window) with
the retraining axis (single fit versus walk-forward refits every 63
snapshots). All four cells used one shared data fetch (all fifty names
this time, 68{,}200 rows), the same split, the same seed, and carried the
two no-learning anchors of \cref{sec:anchors}.
\Cref{tab:eq-ablation} shows the result.

\begin{table}[htbp]
  \centering
  \caption{The graph~$\times$~retraining ablation, one seed, IC loss.
  ``Att.\ VA'' is the attention value-add: GAT OOS Sharpe minus the
  uniform anchor's. Anchors, for reference: island mean Sharpe $-2.13$;
  uniform anchor $-1.05$ (static graph) and $-1.19$ (dynamic).}
  \label{tab:eq-ablation}
  \begin{tabular}{lS[table-format=-1.4,round-precision=4]S[table-format=-1.2,round-precision=2]lS[table-format=-1.2,round-precision=2]}
    \toprule
    Configuration & {OOS IC} & {OOS Sharpe} & Gates & {Att.\ VA} \\
    \midrule
    static + single       & -0.0051 &  1.04 & F/F/T/T & +2.09 \\
    dynamic + single      & -0.0348 & -1.01 & F/F/T/T & +0.18 \\
    static + walk-forward & +0.0200 &  1.95 & F/T/T/T & +3.00 \\
    dynamic + walk-forward& +0.0007 &  0.05 & F/T/T/T & +1.24 \\
    \bottomrule
  \end{tabular}
\end{table}

Two findings, one of them negative. Walk-forward retraining is the lever:
it lifts the static-graph cell from an OOS IC of $-0.005$ to $+0.020$ and
flips the Consistency gate to pass. The decay diagnosed in
\cref{sec:eq-first} was model staleness. The dynamic graph, by contrast,
hurts in both arms --- per-snapshot correlation estimates from a trailing
window are noisy, and the frozen train-period graph evidently acts as a
regulariser. This negative result is worth having: the intuitively
appealing refinement (``keep the graph fresh'') makes things worse as
implemented, and the project kept the static graph with its mild,
disclosed in-sample lookahead rather than trade OOS performance for
cosmetic cleanliness.

The attention A/B, meanwhile, is positive in every cell. Both no-learning
anchors are clearly negative out-of-sample; the GAT beats the uniform
anchor by $+0.2$ to $+3.0$ Sharpe on the same topology with the same
inputs, and its scores correlate with the uniform anchor's at Spearman
$-0.29$ to $-0.14$ --- it is not learning to smooth, it is learning
something else. Against the strict value-added bar the best cell still
falls short (best single at 3.07; the gap narrows from $-2.03$ to
$-1.12$), but the core-thesis comparison holds regardless of which cell
one reads.

One warning surfaced here as well: the static+single cell read
$-0.0051$/1.04 where the near-identical first run had read $+0.0066$/1.42.
Run-to-run variance was clearly material, which is what the next section
is about.

\section{Variance: seeds, devices, and the reporting protocol}
\label{sec:eq-variance}

Five seeds per arm on the same panel put numbers on the spread. The
single-fit arm scored an OOS Sharpe of $0.73 \pm 0.25$ (per-seed range
0.36--1.04) and the walk-forward arm $1.15 \pm 0.75$ (range 0.13--1.95);
walk-forward's mean advantage came with three times the variance, and a
Welch $t$ of about 1.2 at $n=5$ per arm. The single-seed reading of the
ablation --- ``walk-forward is the lever'' --- had ridden the luckiest
walk-forward draw. Meanwhile the attention value-add was positive in all
ten runs, spanning $+1.18$ to $+3.00$ Sharpe against a seed-independent
anchor: the core claim survived the variance study; the walk-forward
claim was downgraded to ``improves the mean, adds variance, not
significant at $n=5$''.

A hardware benchmark added an unplanned second lesson. Moving training to
a laptop GPU was 24--32\% \emph{slower} at this graph size --- fifty
nodes and ten features per snapshot is latency-bound, and kernel-launch
overhead dominates --- so production runs stayed on CPU. More
interesting was that the same seed on a different device swung the result
by the full width of the seed distribution. The seed fixes the initial
weights, but dropout masks come from each device's own random stream and
parallel reductions round in a different order, so training diverges from
the first step; a cross-device rerun is another draw from the same
distribution, not the same experiment on new hardware.

Together these fixed the reporting protocol for everything that follows:
every headline number is a mean $\pm$ standard deviation over at least
five seeds on a fixed device, and what the project claims to be
reproducible is the distribution, not bit-exact point estimates.

\section{Hyperparameter selection with clean hands}
\label{sec:eq-hp}

The hyperparameter search was designed so that the OOS window could not
influence selection even in principle. Twenty-four configurations
(learning rate $\{5\times10^{-4}, 10^{-3}, 3\times10^{-3}\}$, hidden size
$\{32, 64\}$, heads $\{2, 4\}$, layers $\{1, 2\}$) ran three seeds each
--- 72 training runs scored exclusively on the double-embargoed
validation IC, with no OOS metric computed anywhere in the grid. Only the
winner then touched the OOS window, once, with fresh seeds.

The grid's surface is flat on top: six configurations sit within
validation IC 0.10--0.115, and the shipped default ranked third. The one
structural regularity is depth --- all six top configurations have two
GAT layers. Results are therefore robust to reasonable hyperparameter
choices, which is simultaneously a robustness point and a caution against
reading much into tuning.

The winner's OOS validation initially produced a finding that did not
survive its own scrutiny. The validation run had silently executed on the
GPU --- the CUDA wheel had just been installed and the fitting routine
auto-selected it --- while the reference numbers it was compared against
were CPU runs. Given the cross-device divergence of
\cref{sec:eq-variance}, that comparison confounded device with
configuration. The code was fixed the same day (CPU is now an explicit
default, CUDA opt-in) and the validation rerun on CPU. The deconfounded
result: in the single arm the winner is a wash against the default (0.80
vs.\ 0.73 Sharpe, IC near zero in both), but in the walk-forward arm it
is better on the mean and far tighter --- OOS IC $0.0148 \pm 0.0043$
against the default's $0.0055 \pm 0.0147$, a 2.7-fold mean with a third
of the variance, all five seeds positive on both IC and Sharpe, and the
IC mean 3.4 standard deviations above zero. Within this configuration,
walk-forward versus single reaches significance (Welch $t \approx 3.9$ on
IC, $\approx 2.4$ on Sharpe at $n=5$ per arm), upgrading the walk-forward
claim to ``significant in the tuned configuration, directionally
consistent everywhere''.

The best known setup is therefore the static graph, IC loss,
walk-forward refits, and the winner hyperparameters, at an OOS IC of
$0.0148 \pm 0.0043$ and OOS Sharpe of $1.37 \pm 0.39$. Across the three
run families --- seed study, GPU validation, CPU rerun --- the attention
value-add is positive in 30 of 30 seeded runs, spanning $+0.69$
to $+3.03$ Sharpe. That, not any single point estimate, is the equity
track's central result.

\section{What the attention learned}
\label{sec:eq-attention}

The A/B says learned attention beats uniform averaging; attention
analysis on a representative trained model (919{,}336 edge-weight rows
over 1{,}364 dates) says how.

The attention is genuinely relational. The mean self-loop weight is
0.084 --- the model places 91.6\% of its attention mass on neighbours ---
which rules out the trivial failure mode in which attention collapses
onto each node's own features and the graph is decoration. But across a
node's roughly 12.5 neighbours the distribution is nearly uniform:
normalised entropy 0.956, a top-neighbour share of 0.152 where exact
uniformity would give 0.08, and a sector-homophily lift of just $+0.019$
(attention-weighted same-sector share 0.330 against a structural 0.311).
The trained model is a gentle reweighting of mean pooling, not a sharp
selector. This is the mechanism behind the whole chapter: small, learned
tilts on top of broad pooling explain why the GAT's edge over the uniform
anchor is consistent yet modest, and why its scores are nearly
uncorrelated with the anchor's.

The coarse structure is also temporally stationary: the self/neighbour
split holds at 0.91--0.92 and the homophily lift stays within
$+0.01$ to $+0.03$ across the full 2021--2026 span, with no visible break
at the IS/OOS boundary (\cref{fig:att-neighbour}). The project's original design hypothesis --- that
attention would adapt to macro regimes --- is not supported at this
aggregate level, and the report states so; whatever adaptation exists
lives in per-name weights inside a stable coarse structure. Finally, the
names drawing the most incoming attention are mega-caps and sector
bellwethers --- AMZN, NFLX, AVGO, GS, CAT, NVDA, MSFT, LIN --- exactly
the names the rest of the universe co-moves with, a qualitative validity
check the attention passes.

\begin{figure}[htbp]
  \centering
  \includegraphics[width=\textwidth]{att_neighbour_weight.png}
  \caption{Attention mass placed on neighbours (one minus the
  self-loop weight) per snapshot over 2021--2026, from the archived
  attention analysis of the representative trained model. The share
  stays in a narrow band around 0.91--0.92, with no break at the
  dashed IS/OOS boundary --- the stationarity reported in the text.}
  \label{fig:att-neighbour}
\end{figure}

\section{Summary of the equity evidence}
\label{sec:eq-summary}

Learned attention over the correlation graph reliably beats unlearned
propagation over the same graph with the same inputs: positive attention
value-add in 30 of 30 seeded runs across two hyperparameter
configurations and two devices. The resulting composite is a real,
unique, consistent factor that beats nine of its ten inputs but not the
best of them; the value-added gate remains the one honest failure.
Dynamic graphs hurt, walk-forward retraining helps (significantly so in
the tuned configuration), single-seed and single-device numbers are not
to be trusted, and the attention mechanism itself is interpretable ---
relational, nearly uniform, stationary, and centred on economically
sensible hubs.
